\chapter{Limitations and Future Enhancements}
\label{chap:limitations}

\section{System Limitations}
While the deployed Cloud Campus Management System successfully fulfills all functional requirements and operates within the AWS Free Tier, specific technical limitations exist in the current implementation:
\begin{enumerate}
    \item \textbf{AWS Lambda Cold Start Overhead}: Initial API calls following periods of inactivity experience a 240ms–380ms cold start delay as the microVM container is provisioned inside the VPC.
    \item \textbf{Absence of Edge CDN (Amazon CloudFront)}: Static assets are served directly from S3 without an edge content delivery network (CDN), resulting in slight latency variations for geographically distant users.
    \item \textbf{Database Connection Limits on Serverless}: Pairing serverless Lambda with Amazon RDS PostgreSQL requires strict connection pooling management to avoid exhausting database connections during sudden concurrency spikes.
    \item \textbf{Manual Domain Name Assignment}: The application currently uses standard API Gateway and S3 endpoints rather than a custom institution domain configured via Amazon Route 53.
\end{enumerate}

\section{Future Enhancements}
Future development iterations will incorporate the following enterprise cloud architecture extensions:
\begin{itemize}
    \item \textbf{Amazon CloudFront Integration}: Deploying a CloudFront CDN distribution in front of S3 and API Gateway to cache static assets at 450+ global edge locations and lower response latency to under 10ms.
    \item \textbf{AWS WAF (Web Application Firewall)}: Integrating AWS WAF to filter SQL injection, cross-site scripting (XSS), and rate-based DDoS attacks at the API Gateway boundary.
    \item \textbf{RDS Proxy Connection Pooling}: Provisioning Amazon RDS Proxy between AWS Lambda and PostgreSQL to pool and share database connections, enabling support for tens of thousands of concurrent serverless requests.
    \item \textbf{Amazon Route 53 Custom DNS}: Configuring custom SSL/TLS certificates via AWS Certificate Manager (ACM) and mapping institution domain names (\texttt{campus.university.edu}).
    \item \textbf{Automated CI/CD Pipeline}: Establishing an automated continuous integration and continuous deployment (CI/CD) pipeline using AWS CodePipeline and GitHub Actions for automated linting, testing, and cloud deployment upon Git push.
\end{itemize}
